% !TeX program = xelatex
\documentclass[12pt,a4paper]{article}
\usepackage[margin=2cm]{geometry}
\usepackage{fontspec}
\usepackage{amsmath,amssymb}
\setmainfont{Times New Roman}
\setlength{\parindent}{0pt}
\setlength{\parskip}{6pt}
\begin{document}
\begin{center}\Large\bfseries Đề thi vào lớp 10 Hà Nội năm 2013\\Phần Đại số\end{center}

\textbf{Câu 21.} Với $x>0$, cho hai biểu thức
\[
A=\frac{2+\sqrt{x}}{\sqrt{x}},\qquad
B=\frac{\sqrt{x}-1}{\sqrt{x}}+\frac{2\sqrt{x}+1}{x+\sqrt{x}}.
\]
\begin{enumerate}
\item[a)] Tính giá trị của $A$ khi $x=64$.
\item[b)] Rút gọn biểu thức $B$.
\item[c)] Tìm $x$ để $\dfrac AB>\dfrac32$.
\end{enumerate}

\textbf{Câu 22.} Giải bài toán sau bằng cách lập phương trình. Quãng đường từ $A$ đến $B$ dài $90$ km. Một người đi xe máy từ $A$ đến $B$. Khi đến $B$, người đó nghỉ $30$ phút rồi quay về $A$ với vận tốc lớn hơn vận tốc lúc đi $9$ km/h. Tổng thời gian từ lúc bắt đầu đi đến khi trở về $A$ là $5$ giờ. Tính vận tốc xe máy lúc đi từ $A$ đến $B$.

\textbf{Câu 23.}
\begin{enumerate}
\item[1)] Giải hệ phương trình
\[
\begin{cases}
3(x+1)+2(x+2y)=4,\\
4(x+1)-(x+2y)=9.
\end{cases}
\]
\item[2)] Cho parabol $(P):y=\dfrac12x^2$ và đường thẳng $(d):y=mx-\dfrac12m^2+m+1$.
\begin{enumerate}
\item Với $m=1$, xác định tọa độ giao điểm $A,B$ của $(d)$ và $(P)$.
\item Tìm $m$ để $(d)$ cắt $(P)$ tại hai điểm phân biệt có hoành độ $x_1,x_2$ thỏa mãn $|x_1-x_2|=2$.
\end{enumerate}
\end{enumerate}

\textbf{Câu 25.} Với $a,b,c$ là các số dương thỏa mãn
\[
a+b+c+ab+bc+ca=6abc,
\]
chứng minh rằng
\[
\frac1{a^2}+\frac1{b^2}+\frac1{c^2}\ge 3.
\]
\end{document}
